\section*{第 \ref{ch:b3:spin-two-level} 章——自旋与两能级系统}

\begin{solution}{exo:b3:spin-two-level:1}
(a) 直接相乘。(b) $[\hat S_x, \hat S_y] =
(\hbar/2)^2\,2\iu\sigma_z = \iu\hbar\hat S_z$。(c) $\hat S^2 =
(\hbar/2)^2 \times 3\,\mathbb 1 = \tfrac34\hbar^2$，确为
$\tfrac12\cdot\tfrac32\hbar^2$。(d) $\{\mathbb 1, \sigma_x,
\sigma_y, \sigma_z\}$ 张成
厄米 $2\times2$ 矩阵构成的四维实空间（$a + \vect b\cdot\vect\sigma$，其中 $a,
\vect b$ 均为实数）。
\end{solution}

\begin{solution}{exo:b3:spin-two-level:2}
(a) $\ket{\pm x} = (\ket\uparrow \pm \ket\downarrow)/\sqrt2$；
$\ket{\pm y} = (\ket\uparrow \pm \iu\ket\downarrow)/\sqrt2$。(b)
各占一半。(c) 仍各占一半：$|\braket{\pm y}{+x}|^2 =
|1 \mp \iu|^2/4 = \tfrac12$。(d) $\ket{+x}$：赤道上 $\varphi =
0$；沿 $z$ 或 $y$ 方向测量时，投影到相距 $90^\circ$ 的极点上
——总是 $\cos^2(45^\circ) = \tfrac12$。
\end{solution}

\begin{solution}{exo:b3:spin-two-level:3}
(a) $\cos^2(\theta/2)$。(b) $\cos^2(\theta/2) \times
\cos^2(\theta/2) = \cos^4(\theta/2)$。(c) $\theta = 90^\circ$：
$(\tfrac12)^2 = \tfrac14$，即起偏器链的结果。(d)
显然 $\theta = 0$；真正的启示在于\emph{温和}的
中间测量代价最小——这正是
\cref{exo:b3:quantum-formalism:5}的小步极限，许多小转动使态通过时
损耗趋于零。
\end{solution}

\begin{solution}{exo:b3:spin-two-level:4}
(a) $2\mu_{\text{B}}B = \qty{116}{\micro eV}$：$\nu =
\qty{28}{GHz}$。(b) $2 \times 2.79\,\mu_{\text{N}}B =
\qty{0.18}{\micro eV}$：$\nu = \qty{42.6}{MHz}$。(c)
与 $k_{\text{B}}T = \qty{25.9}{meV}$ 相比，电子极化
度为 $\sim\num{2e-3}$，质子 $\sim\num{3e-6}$——热
平衡下的\omterm{def:b3:spin-two-level:spin}{自旋}系综几乎完全无序。(d) ESR：微波；
NMR：射频——两者相距三个数量级，同属一种物理。
\end{solution}

\begin{solution}{exo:b3:spin-two-level:5}
(a) $\dd\langle\hat S_x\rangle/\dd t = -\gamma B_0\langle\hat
S_y\rangle$，$\dd\langle\hat S_y\rangle/\dd t = \gamma
B_0\langle\hat S_x\rangle$，$\dd\langle\hat S_z\rangle/\dd t = 0$。
(b) $\langle\hat S_x\rangle = \tfrac\hbar2\cos\omega_0t$，
$\langle\hat S_y\rangle = \tfrac\hbar2\sin\omega_0t$（符号
由 $\gamma$ 的符号决定）。(c) 两者均由(b)直接看出。(d) 态是
两个能量本征态的\emph{叠加}；其相对
相位以 $\omega_0$ 推进，而这一旋转相位\emph{正是}
进动——定态能级与转矢是
同一两能级演化的两种读法。
\end{solution}

\begin{solution}{exo:b3:spin-two-level:6}
(a) 代入并利用半角公式。(b)
$\cos\tfrac\theta2\cos\tfrac{\pi-\theta}2 +
\eu^{\iu\pi}\sin\tfrac\theta2\sin\tfrac{\pi-\theta}2 =
\cos\tfrac\theta2\sin\tfrac\theta2 -
\sin\tfrac\theta2\cos\tfrac\theta2 = 0$。(c) 赤道上取
$\varphi = 0, \pi$（$\sigma_x$）与 $\varphi = \pm\pi/2$
（$\sigma_y$）处。(d) 绕任一赤道轴的 $180^\circ$ 转动
——\omterm{prop:b3:spin-two-level:rabi}{磁共振}中的 $\pi$ 脉冲即量子比特的 NOT 门。
\end{solution}

\begin{solution}{exo:b3:spin-two-level:7}
(a) 在旋转参考系中驱动场静止，而参考系的
旋转从有效静磁场中减去 $\omega/\gamma$
（与旋转参考系中惯性力的处理方式相同）。(b) 只剩
$B_1$：布洛赫矢量绕（固定的）
横向 $B_1$ 轴以 $\Omega_1 = |\gamma|B_1$ 进动——稳定
翻转。(c) $t_\pi = \pi/\gamma B_1 = 1/(2 \times 42.58\,\text{
MHz/T} \times 10\,\unit{\micro T}) = \qty{1.2}{ms}$。(d) 一条快速
收紧锥面的螺旋：绕 $z$ 轴的进动频率为 \qty{128}{MHz}，
同时以 \qty{426}{Hz} 缓慢向下章动——一条千圈的
螺线，从一极到另一极。
\end{solution}

\begin{solution}{exo:b3:spin-two-level:8}
(a) $\nu = \Delta E/h = \num{5.9e-6}/\num{4.14e-15} =
\qty{1.42e9}{Hz}$；$\lambda = c/\nu = \qty{21.1}{cm}$。(b) $\Delta
\nu \sim 1/2\pi\tau \sim \qty{e-15}{Hz}$：完全可忽略——
每个观测到的线宽都源于运动。(c) $\Delta\nu =
\nu\,(2v/c) = 1420\,\text{MHz} \times \num{1.5e-3} \approx
\qty{2}{MHz}$：旋转圆盘的双峰轮廓。(d)
凡氢呈原子态（温暖、弥散）之处，21 cm 线都可显示；CO
则显示氢已结合为不可见的 H$_2$（寒冷、
致密）：二者合用，可盘点一个星系的整个星际
介质。
\end{solution}

\begin{solution}{exo:b3:spin-two-level:9}
(a) 两种构型，经隧穿耦合：低能时态空间为
二维的——氨\emph{就是}一个
伪装的\omterm{def:b3:spin-two-level:spin}{自旋} $\tfrac12$ 体系。(b) $\nu = \Delta E/h = \qty{24}{GHz}$，$\lambda =
\qty{1.25}{cm}$。(c) 束流到达时，必须使\emph{上}
能级布居过多（由静电偏转选出）——
即粒子数反转，热平衡下不可获得，因为玻尔兹曼
总是偏爱低能级。(d) 受激放大靠
粒子数反转——激光器的工作原理，
首先在微波频率上得到演示。
\end{solution}

\begin{solution}{exo:b3:spin-two-level:10}
(a) $\mu_{\text{B}} \times \qty{10}{T} = \qty{0.58}{meV}$，
与 $\alpha^2E_{\text{I}} \approx \qty{0.7}{meV}$ 量级相当，见
\cref{exo:b3:hydrogen-atom:11}：精细结构\emph{正是}\omterm{def:b3:spin-two-level:spin}{自旋}
与运动磁场相遇。(b) $\Delta E = hc\,\Delta\lambda/
\lambda^2 = \qty{2.1}{meV}$：内部场 $\Delta E/2\mu_{
\text{B}} \approx \qty{18}{T}$。(c) 内场随
$Z^4$ 增长（核电荷的三次方在场中，再一次在轨道
半径中）：重原子的精细结构用袖珍
分光镜直接看到。(d) $j = \tfrac32$：四个态；$j = \tfrac12$：两个
——共六个，恰为 $2 \times (2l + 1)$，当 $l = 1$ 时。
\end{solution}

\begin{solution}{exo:b3:spin-two-level:11}
(a) $S_z$ 的取值为 $+\hbar, 0, 0, -\hbar$；$\hat S^2$ 作用于
对称组合上（借助\omterm{def:b3:harmonic-oscillator:ladder}{升降算符}恒等式）得 $2\hbar^2$（$S = 1$），
反对称组合上得 $0$（$S = 0$）。(b)
只有单态（以及中间的三重态）不能写成
乘积态——单态才是\emph{真正的}最大纠缠
对。(c) 发射意味着三重态位于上方：平行\omterm{def:b3:spin-two-level:spin}{自旋}组态
能量高出 \qty{5.9}{\micro eV}。(d)
每位观察者单独看到的结果完全随机；
关联只有当两份记录\emph{合在
一起}时才出现——局域统计并未改变，因此
没有任何信息被传递（论证见\cref{rem:b3:quantum-formalism:nosignal}）。
\end{solution}

\begin{solution}{exo:b3:spin-two-level:12}
(a) 第一个 $\pi/2$ 脉冲把布洛赫矢量置于赤道面上；在
$T$ 时间内，它以原子自身的 $\nu_0$ 进动，
本地振荡器以 $\nu$ 转动：累积的角度差为
$2\pi\delta T$；第二个 $\pi/2$ 脉冲把该相位转化为
布居数——时间上的干涉仪。(b) $\Delta\delta \sim
1/2T = \qty{1}{Hz}$. (c) $\qty{e-4}{Hz}/\qty{9.19e9}{Hz} \approx
\num{e-14}$——每三百万年差一秒。(d) 同样的绝对条纹分裂
落在高 $10^5$ 倍的载波上，使相对精度提高
$10^5$ 倍：因此锶和镱光学钟的精度达
$10^{-18}$，并且秒的
重新定义即将进行。
\end{solution}

\begin{solution}{pb:b3:spin-two-level:1}
\textbf{1.} $\nu = 42.58 \times 3.0 = \qty{127.7}{MHz}$; $h\nu =
\qty{5.3e-7}{eV}$.
\textbf{2.} $\Delta n/n = h\nu/2k_{\text{B}}T = \num{5.3e-7}/(2
\times 0.0267) \approx \num{e-5}$：百万分之十。
\textbf{3.} $\num{6.6e22}\,\unit{cm^{-3}} \times \num{e-5} \approx
\num{7e17}$ 个携带信号的质子（每立方厘米）——单个\omterm{def:b3:spin-two-level:spin}{自旋}微弱，
但数量惊人。
\textbf{4.} 施特恩–盖拉赫实验逐个测量原子，并按
\omterm{def:b3:quantum-formalism:observable}{本征值}将其分开；MRI 聆听的是 $10^{23}$ 个\omterm{def:b3:spin-two-level:spin}{自旋}的\emph{热平均}，
玻尔兹曼使该平均值保持在距零微伏量级之内。
\textbf{5.} 极化 $\propto B_0$，感应电动势
$\propto \omega \propto B_0$：信号大致 $\propto B_0^2$——
这正是 \qty{3}{T} 优于 \qty{1.5}{T} 的理由，也适用于
\qty{7}{T} 科研设备。
\textbf{6.} 处于杂散梯度中的铁磁体受到
$F = \mu\,\partial B/\partial z$ 的力，按千克量级的铁换算：
数百牛顿的力出现在门口——这正是
设置屏蔽措施与检查清单的原因。
\textbf{7.} $\nu_1 = \gamma B_1/2\pi = \qty{426}{Hz}$: $\pi/2$
脉冲 \qty{0.59}{ms}，$\pi$ 脉冲 \qty{1.2}{ms}。
\textbf{8.} 它在横向平面内以
\qty{127.7}{MHz} 进动；旋转磁化强度穿过
线圈的通量按法拉第定律感应出射频电动势，频率恰为该值——
即原始 MR 信号。
\textbf{9.} 每一次交换能量的翻转也会使翻转\omterm{def:b3:spin-two-level:spin}{自旋}的
相位随机化：任何导致 $T_1$ 的因素也会贡献于 $T_2$，
而退相位还有其自身的额外通道——相干性
最先消亡。
\textbf{10.} 在 $\tau$ 时刻，每位跑者都折返：跑得快的，
此时领先最远，回程也最远；在 $2\tau$ 时刻，所有人
并排回到起跑线——退相位得以倒转，并
回波随之响起。
\textbf{11.} 只有\emph{静态}部分可被倒转：一个保持
恒定（即使有误差）频率的\omterm{def:b3:spin-two-level:spin}{自旋}，其相位被精确回溯；
而涨落的分子场在两半之间会发生变化，因此
无法倒转——它们的衰减才是真正的、组织特异的 $T_2$。
\textbf{12.} 早且频繁地采样（短 TR、短 TE），
脂肪的短 $T_1$ 格外醒目；若等待很久、回波很晚，长 $T_2$
的流体会发亮：序列的时钟设置决定了由哪个弛豫常数
来绘制图像。
\textbf{13.} $\dd\nu/\dd z = 42.58\,\unit{MHz/T} \times
0.04\,\unit{T/m} = \qty{1.7}{kHz/mm}$.
\textbf{14.} 局部拉莫尔频率落入该频带内的薄层，
即厚度 $\qty{2}{kHz}/\qty{1.7}{kHz/mm} \approx
\qty{1.2}{mm}$——即选定的切片。
\textbf{15.} 傅里叶变换：频率标记位置，因此
回波的频谱\emph{正是}该切片的投影。
\textbf{16.} $256^2 \times 12 \approx \qty{100}{kB}$；每个
重复时间（量级为 $T_1$）编码一行，$256$ 行需
数分钟——这正是要求患者保持静止的原因。
\textbf{17.} 梯度线圈在切换时承载千安量级的开关电流，
置于 \qty{3}{T} 磁场中，拉普拉斯力在每次切换时都
将其撞向支架——这正是
每次扫描中机关枪般声响的来源。
\textbf{18.} X 射线反映电子密度——骨与空气的差别——
并沉积电离剂量；MRI 读取质子密度与两种
弛豫时钟，它们在软组织间差异丰富：脑、
软骨与肿瘤，对 X 射线而言毫无二致，对\omterm{def:b3:spin-two-level:spin}{自旋}而言却各不相同，
且电离剂量为零。
\textbf{19.} 每种化学环境使其质子屏蔽约
百万分之几，分子的质子呈现为分辨清晰的
位移谱线梳——试管中的结构测定。
\textbf{20.} 活动提高局部血流量与氧合水平；
顺磁性的脱氧血红蛋白被稀释；局部 $T_2^*$ 延长；
体素在思维之后数秒即变亮。
\textbf{21.} 在 \qty{84}{GHz} 下，组织不透明（毫米
波几乎无法穿透皮肤），且电子弛豫
仅为微秒量级——活体中无望，但非常适合计数
材料中的自由基与缺陷，以及剂量测定，堪称完美。
\textbf{22.} $\num{5e-7}\,\unit{eV}$ 与 eV 量级的化学键相比：
弱了约千万倍，不足以打断任何化学键——非电离；一个
\qty{60}{keV} 的 X 射线光子携带的能量则多 $10^{11}$ 倍。
\textbf{23.} 从 $10^{-5}$ 提升到量级为一：每个
原子核增益 $\sim10^5$——这是必要的，因为吸入气体比水中的质子稀薄
数千倍。
\textbf{24.} \omterm{def:b3:spin-two-level:spin}{自旋} $\tfrac12$ 体系是\emph{完美的}两能级
系统（无泄漏能级），它对电荷噪声的耦合
仅经由磁矩而微弱，因此在工业材料中
可获得长相干性。
\textbf{25.} 一个 \qty{128}{MHz} 进动，一个 $10^{-5}$
极化、两个弛豫时钟与三个梯度：\omterm{def:b3:spin-two-level:spin}{自旋}
物理学逐层为活体大脑拍照，
电离剂量为零——1922 年的两个银斑已长成
一个医院科室。
\end{solution}
